% TOP_PROBLEM: {"id": 1319, "title": "1/3–2/3 Conjecture", "category_id": 2, "category": {"id": 2, "name": "combinatorics", "display_name": "Combinatorics", "description": "Counting problems, graph theory, discrete structures.", "slug": "combinatorics", "order_index": 2, "created_at": "2026-07-31T15:26:25.671Z"}, "status": "open"}
% FIRST_POSTED: 2026-10-01
% ATTEMPT_STATUS: unresolved
% Imported from: unsolvedmath_additional_500.tex; UnsolvedMath ID 1319
% !TeX program = lualatex
% Compile twice: lualatex 1319.tex
% Self-contained: no external bibliography, images, or \input files.
% Numeric section labels and visible numbers are original UnsolvedMath IDs.
\documentclass[11pt,a4paper]{article}
\usepackage[margin=24mm,headheight=14pt]{geometry}
\usepackage{fontspec}
\setmainfont{Latin Modern Roman}
\setsansfont{Latin Modern Sans}
\setmonofont{Latin Modern Mono}
\usepackage{amsmath,amssymb,mathtools}
\usepackage{unicode-math}
\setmathfont{Latin Modern Math}
\usepackage{microtype}
\usepackage{enumitem,xurl,needspace}
\usepackage[dvipsnames]{xcolor}
\usepackage{hyperref}
\hypersetup{unicode=true,colorlinks=true,linkcolor=MidnightBlue,urlcolor=MidnightBlue,
 pdftitle={UnsolvedMath Problem 1319},pdfauthor={Source-annotated compilation},bookmarksnumbered=true}
\usepackage{bookmark,tocloft,titlesec,fancyhdr}
\setlength{\cftsecnumwidth}{4.2em}
\cftsetpnumwidth{2.5em}
\cftsetrmarg{4em}
\titleformat{\section}{\Large\bfseries\raggedright}{\thesection}{0.7em}{}
\pagestyle{fancy}\fancyhf{}
\fancyhead[L]{\small\sffamily Additional UnsolvedMath statements}
\fancyhead[R]{\small Original catalogue IDs}
\fancyfoot[C]{\thepage}
\setlength{\parindent}{0pt}
\setlength{\parskip}{5pt plus 1pt}
\setlength{\emergencystretch}{3em}
\setcounter{tocdepth}{1}\setcounter{secnumdepth}{1}
\setlist[itemize]{leftmargin=3em,itemsep=4pt,topsep=4pt,parsep=0pt}
\allowdisplaybreaks
\newcommand{\N}{\mathbb{N}}
\newcommand{\Z}{\mathbb{Z}}
\newcommand{\Q}{\mathbb{Q}}
\newcommand{\R}{\mathbb{R}}
\newcommand{\C}{\mathbb{C}}
\newcommand{\F}{\mathbb{F}}
\newcommand{\A}{\mathbb{A}}
\newcommand{\PP}{\mathbb{P}}
\newcommand{\E}{\mathbb{E}}
\newcommand{\eps}{\varepsilon}
\newcommand{\dd}{\,\mathrm d}
\newcommand{\sref}[2]{\hyperlink{source:#1:#2}{[S#2]}}
\newif\ifshowcatalogue
\showcataloguetrue
\newcommand{\cataloguescope}[1]{\ifshowcatalogue
 \paragraph{Statement recorded in the supplied source.}
 #1\fi}
\begin{document}
% UnsolvedMath ID 1319; source catalogue code COMB-001

\setcounter{section}{1318}

\section{The One-Third--Two-Thirds Poset Conjecture}\label{1319}

{\small\sffamily UnsolvedMath ID 1319 · Combinatorics}

\subsection{Definitions and mathematical statement}

\paragraph{Shared notation.}Unless a section says otherwise, $\N=\{1,2,\ldots\}$,
$\Z$ is the ring of integers, $\Q,\R,\C$ are the rational, real, and complex
fields, $[N]=\{1,\ldots,\lfloor N\rfloor\}$, and $\log$ is the natural
logarithm. For real functions of a parameter tending to infinity,
$f\ll g$ and $f=O(g)$ mean $|f|\leq Cg$ eventually for some fixed $C$;
$f\gg g$ reverses this comparison; $f=o(g)$ means $f/g\to0$;
$f\sim g$ means $f/g\to1$; and $f\asymp g$ means both order bounds.
A subscript on $O$ or $\ll$ specifies allowed constant dependence.
The expression $n^{o(1)}$ in an upper bound means growth smaller than
$n^\epsilon$ for each fixed $\epsilon>0$. Local definitions govern any
exception, finite-field convention, or use of the same symbol for another
object. An asymptotic claim is not automatically an effective algorithm.



A linear extension of a finite partially ordered set $P$ is a total order on its elements respecting all comparisons in $P$. Choose an extension uniformly among the finitely many possibilities. The conjecture seeks distinct elements $x,y$ with $1/3\leq\Pr(x\text{ precedes }y)\leq2/3$, provided $P$ is not itself totally ordered. \sref{1319}{1} \sref{1319}{2}

\paragraph{Statement recorded in the supplied source.}
Does every non-totally-ordered finite poset have two elements with probability between 1/3 and 2/3 in random linear extensions? \sref{1319}{1}

\paragraph{Residual task reported by the archive.}

The embedded review (2026-08-17) records: Establish a balanced incomparable pair in every finite non-total poset. \sref{1319}{1}

\subsection{Short English statement}

Must a finite partial order that is not already a total order contain a pair whose relative order is reasonably balanced across all linear extensions?

\subsection{Research attempt: exact extension counts and symmetry witnesses}
\textbf{Status and dated context.} No proof for all finite posets is obtained. The inspected primary abstract [R1] proves balanced pairs in several structured classes, including the automorphism phenomenon used below. A separate primary preprint posted 25 September 2026 [R2] reports an improvement in a general lower bound, not the full one-third bound. We do not infer a complete solution from either abstract.

\subsubsection{1. A finite exact recurrence}
Represent a poset on $n$ elements by predecessor sets $D(x)$. For an order ideal $I$ let $F(I)$ count linear extensions of its induced order. Then
\[F(\varnothing)=1,\qquad F(I)=\sum_{x\in\max I}F(I\setminus\{x\}).\tag{1}\]
Every extension has a unique last element in $\max I$, and deleting it gives the indicated bijection. Equivalently count extensions by successively choosing elements all of whose predecessors have already appeared. This uses integers throughout and computes $e(P)=F(P)$ in at most $2^n$ subset states, with polynomial overhead per state.

For incomparable $x,y$, add the relation $x<y$ and take transitive closure, obtaining $P_{x<y}$. Its extensions are exactly those extensions of $P$ placing $x$ before $y$. Thus the conjecture for this pair is the exact integer comparison
\[e(P)\leq3e(P_{x<y})\leq2e(P).\tag{2}\]
This supplies an exhaustive certificate for any finite input, but the number of inputs is unbounded; enumerating them does not prove the universal assertion.

\subsubsection{2. A swapping automorphism forces balance one half}
Suppose an automorphism $\tau$ interchanges $x,y$ (it may permute other elements). Applying $\tau$ to every entry of an extension is a bijection of the extension set. It maps extensions with $x$ before $y$ to those with $y$ before $x$, so their counts agree. The pair must be incomparable, since an order automorphism cannot reverse a strict comparison. Hence its probability is exactly $1/2$.

As a concrete sufficient condition, incomparable elements with identical strict predecessor sets and identical strict successor sets can be transposed while all others are fixed. This proves the claim for such twins. Symmetry, rather than independence of all comparisons, is what justifies the equality.

\subsubsection{3. A nonsymmetric family and two closure operations}
Consider a chain $x_1<\cdots<x_k$ and one isolated point $y$. Its $k+1$ linear extensions are the possible insertion positions of $y$, equally likely. Exactly $j$ put $y$ before $x_j$, giving probability $j/(k+1)$. For $k\geq1$, an integer $j$ nearest $(k+1)/2$ gives a number in $[1/3,2/3]$ (for $k=1$ it is $1/2$). This includes examples with no twin pair inside the long chain.

If a component $P$ has a balanced pair, that pair remains balanced in both the ordinal sum $P\oplus Q$ and the disjoint union $P\sqcup Q$. In the ordinal sum every pair of component extensions has exactly one concatenation. In the disjoint union it has exactly $\binom{|P|+|Q|}{|P|}$ interleavings, independent of the chosen extensions. Therefore restriction to $P$ is uniform in either case. This justifies propagation, but does not prove the conjecture for a component that has no balanced pair already known.

\subsubsection{4. Adversarial verification and remaining gap}
The diagnostic enumerates every transitively closed strict relation contained in the natural order on at most five elements, evaluates (1), and checks (2) for each non-chain. Every finite poset admits a labelling by one of its linear extensions, so this enumeration covers every isomorphism type at those sizes, with possible repeats. It also checks the chain-plus-isolated formula exactly. Using one random extension would not supply the uniform counts required in (2), and rounding probabilities near $1/3$ is avoided entirely. What is absent is a structural argument producing a witness in every nonsymmetric indecomposable poset, or an induction that guarantees such a witness survives all required reductions.

\subsection{Sources}

{\small

\begin{itemize}

\item[\hypertarget{source:1319:1}{S1}] ULAM.AI, \emph{UnsolvedMath}, supplied \texttt{problems.json}, record 1319 (COMB-001), statement and background. The embedded literature-review date is 2026-08-17; that is the archive’s review date, not a fresh certification. Dataset landing page: \url{https://huggingface.co/datasets/ulamai/UnsolvedMath}.

\item[\hypertarget{source:1319:2}{S2}] K. P. Bogart and J. P. Trotter, The 1/3--2/3 Conjecture for 5-Thin Posets, SIAM J. Discrete Math. 21 (2007). \url{https://epubs.siam.org/doi/10.1137/0405037}. Bibliographic information imported from the supplied record.

\item[\hypertarget{source:1319:3}{S3}] 1/3--2/3 conjecture overview (). \url{https://en.wikipedia.org/wiki/1/3%E2%80%932/3_conjecture}. Bibliographic information imported from the supplied record.

\item[R1] E. J. Olson and B. E. Sagan, On the 1/3--2/3 Conjecture (2017), inspected abstract.
\url{https://arxiv.org/abs/1706.04985}.
\item[R2] M. Aires, S. H. Chan, I. Pak and G. Panova, Breaking the Infinite Barrier in the 1/3--2/3 Conjecture, posted 25 September 2026; inspected abstract only.
\url{https://arxiv.org/abs/2609.30888}.
\end{itemize}

}

\label{end:1319}
\end{document}
